
\chapter{Expected Results and Evaluation Metrics}
\section{Expected Results and Evaluation Metrics}

\subsection{Expected Results}

The proposed MotionBridge system is expected to provide a low-cost and real-time framework for human arm motion capture, simulation replication, and robotic dataset generation using only standard RGB cameras and open-source software tools. By integrating MediaPipe-based landmark detection, multi-camera triangulation, and PyBullet simulation, the system is expected to generate accurate three-dimensional representations of human arm and hand motion without requiring specialized depth sensors or wearable devices.

The system is expected to successfully detect shoulder, elbow, wrist, and finger landmarks from multiple camera views and reconstruct their corresponding 3D world coordinates through confidence-weighted triangulation. The reconstructed arm motion is expected to be replicated in real time within a PyBullet simulation environment using inverse kinematics-based robotic arm control.

Simultaneously, the system is expected to generate a structured dataset containing timestamped 3D coordinates of arm and finger joints, along with corresponding robotic arm joint angles. This dataset is intended to support future applications in robotic imitation learning, teleoperation, rehabilitation engineering, and human--robot collaboration.

Furthermore, the multi-camera architecture is expected to improve robustness against partial occlusions by utilizing redundant viewpoints, thereby maintaining reliable landmark reconstruction even when one camera view is temporarily obstructed. The complete pipeline is expected to operate on standard CPU-based hardware while maintaining near real-time performance.

\subsection{Evaluation Metrics}

The performance of the proposed system will be evaluated using quantitative and qualitative metrics that assess landmark detection quality, three-dimensional reconstruction accuracy, real-time performance, simulation responsiveness, dataset quality, and robustness to occlusions.

\subsubsection{Landmark Detection Performance}

The effectiveness of MediaPipe Pose and MediaPipe Hands in detecting arm and hand landmarks will be evaluated using standard classification metrics based on correctly and incorrectly detected landmarks.

\paragraph{Precision}

Precision measures the proportion of detected landmarks that are correctly identified.

\begin{equation}
Precision = \frac{TP}{TP+FP}
\end{equation}

where:
\begin{itemize}
    \item TP = True Positives
    \item FP = False Positives
\end{itemize}

\paragraph{Recall}

Recall measures the proportion of actual landmarks that are successfully detected.

\begin{equation}
Recall = \frac{TP}{TP+FN}
\end{equation}

where:
\begin{itemize}
    \item FN = False Negatives
\end{itemize}

\paragraph{F1-Score}

The F1-score provides a balanced measure of precision and recall.

\begin{equation}
F1 = 2 \times \frac{Precision \times Recall}{Precision + Recall}
\end{equation}

Higher precision, recall, and F1-score values indicate better landmark detection performance.

\subsubsection{Three-Dimensional Reconstruction Accuracy}

The accuracy of the reconstructed 3D joint coordinates will be evaluated by comparing reconstructed positions with known reference measurements or predefined calibration poses.

The Mean Reconstruction Error will be calculated as:

\begin{equation}
E_{mean} = \frac{1}{N}\sum_{i=1}^{N}\left\|X_i-\hat{X}_i\right\|
\end{equation}

where:
\begin{itemize}
    \item $X_i$ represents the reference joint position,
    \item $\hat{X}_i$ represents the reconstructed joint position,
    \item $N$ represents the total number of evaluated joints.
\end{itemize}

Lower reconstruction error indicates better triangulation accuracy and more reliable 3D motion capture.

\subsubsection{Real-Time Processing Performance}

The computational efficiency of the proposed system will be evaluated using the following metrics:

\begin{itemize}
    \item Frames Per Second (FPS)
    \item Average processing time per frame
    \item CPU utilization
\end{itemize}

Frames Per Second (FPS) will be computed as:

\begin{equation}
FPS = \frac{\text{Total Processed Frames}}{\text{Total Execution Time}}
\end{equation}

The system is expected to maintain near real-time operation while processing multiple camera streams simultaneously.

\subsubsection{Simulation Replication Latency}

The responsiveness of the robotic arm simulation will be measured by calculating the delay between human motion capture and corresponding robotic arm movement within PyBullet.

Latency will be calculated as:

\begin{equation}
Latency = T_{simulation} - T_{capture}
\end{equation}

where:
\begin{itemize}
    \item $T_{capture}$ is the timestamp of landmark acquisition,
    \item $T_{simulation}$ is the timestamp of simulation update.
\end{itemize}

Lower latency values indicate better real-time motion replication performance.

\subsubsection{Dataset Completeness}

The quality and reliability of the generated dataset will be evaluated by measuring the percentage of frames containing valid joint information.

Dataset Completeness will be calculated as:

\begin{equation}
Completeness = \frac{\text{Valid Frames}}{\text{Total Frames}} \times 100
\end{equation}

A higher completeness percentage indicates successful landmark tracking and dataset generation throughout the recording session.

\subsubsection{Occlusion Robustness}

To evaluate the robustness of the proposed multi-camera framework, controlled occlusion experiments will be performed by partially blocking one or more camera views.

The following factors will be analyzed:

\begin{itemize}
    \item Landmark detection continuity
    \item Reconstruction accuracy under occlusion
    \item Percentage of successfully reconstructed joints
    \item Simulation stability during occluded conditions
\end{itemize}

The multi-camera configuration is expected to reduce tracking failures and improve reconstruction reliability compared to monocular camera systems.

\subsubsection{Overall System Evaluation}

The overall effectiveness of MotionBridge will be assessed by combining the above metrics to determine whether the system can provide accurate, real-time human arm motion capture and robotic dataset generation using only low-cost RGB cameras and standard computing hardware.

Successful completion of the project will be demonstrated through:

\begin{itemize}
    \item Reliable landmark detection,
    \item Accurate 3D joint reconstruction,
    \item Real-time robotic arm simulation,
    \item Low-latency operation,
    \item High dataset completeness,
    \item Robust performance under moderate occlusion conditions.
\end{itemize}

The results obtained from these evaluations will be used to assess the suitability of the proposed framework for future robotic imitation learning and human--robot interaction applications.